\documentclass[aspectratio=169]{beamer}

\usetheme{metropolis}
\usepackage[utf8]{inputenc}
\usepackage[T1]{fontenc}
\usepackage[french]{babel}
\usepackage{tikz}
\usetikzlibrary{shapes.geometric, arrows.meta, positioning}
\usepackage{amsmath}

% ── Palette de couleurs ───────────────────────────────────────────────────────
\definecolor{sage}{RGB}{90, 120, 75}
\definecolor{sagelight}{RGB}{163, 188, 146}
\definecolor{lavender}{RGB}{110, 95, 140}
\definecolor{cream}{RGB}{250, 247, 240}
\definecolor{charcoal}{RGB}{42, 39, 37}
\definecolor{sand}{RGB}{196, 160, 122}
\definecolor{sagesofter}{RGB}{118, 150, 100}

\setbeamercolor{normal text}{fg=charcoal, bg=cream}
\setbeamercolor{alerted text}{fg=lavender}
\setbeamercolor{example text}{fg=sage}
% Frametitle sobre : fond crème, texte sauge
\setbeamercolor{frametitle}{bg=cream, fg=sage}
\setbeamercolor{progress bar}{fg=sand!80, bg=sand!20}
\setbeamercolor{block title}{bg=sage!15, fg=sage}
\setbeamercolor{block body}{bg=sage!6, fg=charcoal}
\setbeamercolor{block title alerted}{bg=lavender!15, fg=lavender}
\setbeamercolor{block body alerted}{bg=lavender!6, fg=charcoal}
\setbeamercolor{block title example}{bg=sagelight!30, fg=sage}
\setbeamercolor{block body example}{bg=sagelight!10, fg=charcoal}
% Palette primaire (page titre + barre du bas) : crème chaude
\setbeamercolor{palette primary}{bg=cream, fg=charcoal}
\setbeamercolor{title separator}{fg=sand}

\metroset{block=fill, progressbar=foot, titleformat=regular}

% ── Page de titre élégante ────────────────────────────────────────────────────
\setbeamertemplate{title page}{
  \begin{tikzpicture}[remember picture, overlay]
    % Fond crème uniforme (déjà défini dans normal text bg)
    % Bande décorative gauche en sauge doux
    \fill[sage!70] (current page.north west) rectangle ++(.22\paperwidth, -\paperheight);
    % Ligne fine sable en bas
    \fill[sand!60] (current page.south west) rectangle ++(\paperwidth, .018\paperheight);
  \end{tikzpicture}
  \vfill
  \begin{beamercolorbox}[wd=.74\paperwidth, sep=1.2em, rightskip=2em, leftskip=.24\paperwidth]{normal text}
    {\usebeamerfont{title}\color{sage}\Large\textbf{\inserttitle}}\\[0.5em]
    {\usebeamerfont{subtitle}\color{charcoal!70}\insertsubtitle}\\[1.2em]
    \textcolor{sand}{\rule{4cm}{0.5pt}}\\[0.8em]
    {\small\color{charcoal!60}\insertauthor}\\[0.2em]
    {\footnotesize\color{charcoal!50}\insertinstitute}\\[0.2em]
    {\footnotesize\color{charcoal!50}\insertdate}
  \end{beamercolorbox}
  \vfill
}

% ── Métadonnées ───────────────────────────────────────────────────────────────
\title{Persona}
\subtitle{Test de personnalité adaptatif par IA}
\author{Yousra Yasmine Henni Mansour \and Samira Fawaz}
\institute{UE TPALT --- Sorbonne Sciences}
\date{Avril 2026}


\begin{document}

% ── Diapo 1 : Titre ──────────────────────────────────────────────────────────
\maketitle

% ── Diapo 2 : Sommaire ───────────────────────────────────────────────────────
\begin{frame}{Sommaire}
  \setbeamertemplate{section in toc}[sections numbered]
  \tableofcontents[hideallsubsections]
\end{frame}


% =============================================================================
\section{Contexte et objectif}
% =============================================================================

\begin{frame}{Qu'est-ce que Persona ?}

  \begin{block}{Objectif}
    Application web de test de personnalité \textbf{adaptatif}, fondé sur le modèle \textbf{Big Five (OCEAN)}.
  \end{block}

  \vspace{1em}

  \begin{columns}[T]
    \column{0.18\textwidth}
    \begin{block}{\centering O}
      \centering\small Ouverture
    \end{block}
    \column{0.18\textwidth}
    \begin{block}{\centering C}
      \centering\small Conscienciosité
    \end{block}
    \column{0.18\textwidth}
    \begin{block}{\centering E}
      \centering\small Extraversion
    \end{block}
    \column{0.18\textwidth}
    \begin{block}{\centering A}
      \centering\small Agréabilité
    \end{block}
    \column{0.18\textwidth}
    \begin{block}{\centering N}
      \centering\small Névrosisme
    \end{block}
  \end{columns}

  \vspace{1em}

  Ce qui distingue Persona des questionnaires classiques :

  \medskip
  \begin{columns}[T]
    \column{0.3\textwidth}
    \textbf{1. Adaptativité}\\
    \small Questions choisies dynamiquement selon les réponses
    \column{0.3\textwidth}
    \textbf{2. Scoring}\\
    \small Scores normalisés 0--100 par dimension
    \column{0.3\textwidth}
    \textbf{3. Rapport IA}\\
    \small Texte personnalisé généré par un LLM
  \end{columns}

\end{frame}


% =============================================================================
\section{Architecture technique}
% =============================================================================

\begin{frame}{Architecture}

  \begin{center}
  \begin{tikzpicture}[
    node distance=0.8cm and 2cm,
    box/.style={
      draw=sage, rounded corners=6pt,
      minimum width=2.6cm, minimum height=1cm,
      align=center, font=\small\bfseries,
      fill=sagelight!25, line width=1.2pt
    },
    ext/.style={
      draw=lavender, rounded corners=6pt,
      minimum width=2.6cm, minimum height=1cm,
      align=center, font=\small,
      fill=lavender!12, line width=1.2pt
    },
    arrow/.style={-{Stealth[length=5pt]}, line width=1.2pt, sage},
    darrow/.style={{Stealth[length=5pt]}-{Stealth[length=5pt]}, line width=1.2pt, charcoal!50},
  ]
    \node[box] (ng)      {Angular 21\\Frontend};
    \node[box, right=2.2cm of ng] (api) {FastAPI\\Backend};
    \node[ext, above right=0.5cm and 1.8cm of api] (claude) {API Claude\\Anthropic};
    \node[ext, below right=0.5cm and 1.8cm of api] (smtp)   {SMTP\\E-mail};
    \node[ext, below=1.1cm of api] (json) {Stockage\\JSON};

    \draw[darrow] (ng)  -- node[above, font=\scriptsize\color{charcoal}]{REST / HTTP} (api);
    \draw[arrow]  (api) -- node[above, font=\scriptsize, sloped]{rapport} (claude);
    \draw[arrow]  (api) -- node[below, font=\scriptsize, sloped]{envoi} (smtp);
    \draw[arrow]  (api) -- node[right, font=\scriptsize]{sessions} (json);
  \end{tikzpicture}
  \end{center}

\end{frame}


% =============================================================================
\section{L'IA dans Persona}
% =============================================================================

\begin{frame}{Deux niveaux d'intelligence artificielle}

  \begin{columns}[T]
    \column{0.48\textwidth}
    \begin{block}{IA (1) — Sélection adaptative}
      \textbf{Moteur à règles (IA symbolique)}
      \begin{itemize}
        \item Analyse les scores partiels après chaque réponse
        \item Sélectionne la dimension \textbf{la moins explorée}
        \item 30 questions annotées $\rightarrow$ 15 posées
        \item Sans LLM : instantané, coût nul
      \end{itemize}
    \end{block}

    \column{0.48\textwidth}
    \begin{block}{IA (2) — Génération du rapport}
      \textbf{LLM : \texttt{claude-opus-4-6}}
      \begin{itemize}
        \item Appelé en fin de test
        \item Reçoit scores + archétype identifié
        \item Produit un rapport en langage naturel
        \item Format JSON imposé par gabarit strict
      \end{itemize}
    \end{block}
  \end{columns}

  \vspace{0.8em}

  \begin{center}
  \begin{tikzpicture}[
    scale=0.72, transform shape,
    node distance=0.3cm and 0.4cm,
    b/.style={draw=charcoal!40, rounded corners=4pt, minimum width=2.2cm,
              minimum height=0.65cm, align=center, font=\scriptsize, fill=cream},
    ai/.style={draw=sage, rounded corners=4pt, minimum width=2.2cm,
               minimum height=0.65cm, align=center, font=\scriptsize\bfseries,
               fill=sagelight!40, line width=1.2pt},
    arr/.style={-{Stealth[length=4pt]}, line width=1pt, sage},
  ]
    \node[b]  (a) {Accueil};
    \node[b,  right=of a] (q) {Question $n$};
    \node[ai, right=of q] (ia1) {IA (1)\\Sélection};
    \node[b,  right=of ia1] (fin) {15 questions};
    \node[ai, right=of fin] (ia2) {IA (2)\\Rapport};
    \node[b,  right=of ia2] (mail) {Résultat\\+ e-mail};

    \draw[arr] (a) -- (q);
    \draw[arr] (q) -- (ia1);
    \draw[arr] (ia1) -- node[above, font=\tiny]{$\times 15$} (fin);
    \draw[arr] (fin) -- (ia2);
    \draw[arr] (ia2) -- (mail);
  \end{tikzpicture}
  \end{center}

\end{frame}


% =============================================================================
\section{Fonctionnalités}
% =============================================================================

\begin{frame}{Scoring et archétypes}

  \begin{columns}[T]
    \column{0.48\textwidth}
    \begin{block}{Normalisation des scores}
      \[
        s_t = \frac{\overline{c_t} + 5}{10} \times 100
        \qquad
        \overline{c_t} = \frac{1}{n_t}\sum_{i=1}^{n_t} a_i \cdot p_i
      \]
      \smallskip
      \small $a_i \in [1,5]$ : réponse \quad $p_i \in \{-1,+1\}$ : polarité\\
      Résultat : $s_t \in [0, 100]$ par dimension
    \end{block}

    \vspace{0.5em}
    \begin{block}{Sélection de l'archétype}
      \[
        d(\text{arch}) = \sqrt{\sum_{t}(s_t - s_t^{\text{arch}})^2}
      \]
      \small Archétype le plus proche du profil mesuré
    \end{block}

    \column{0.48\textwidth}
    \begin{exampleblock}{Les 6 archétypes}
      \begin{itemize}
        \item Le Visionnaire
        \item Le Logicien
        \item Le Gardien
        \item L'Architecte
        \item L'Empathique
        \item L'Explorateur
      \end{itemize}
    \end{exampleblock}
  \end{columns}

\end{frame}

% -----------------------------------------------------------------------------
\begin{frame}{Rapport et envoi par e-mail}

  \begin{columns}[T]
    \column{0.55\textwidth}
    \begin{block}{Contenu du rapport généré par l'IA}
      \begin{itemize}
        \item Archétype + tagline personnalisée
        \item Résumé global du profil
        \item Score par trait (0--100) + interprétation
        \item Points forts
        \item Axes de développement
        \item Recommandations personnalisées
        \item Disclaimer (non médical)
      \end{itemize}
    \end{block}

    \column{0.42\textwidth}
    \begin{block}{Envoi automatique}
      \begin{itemize}
        \item Format HTML
        \item SMTP configurable
        \item Renvoi possible (max 3)
      \end{itemize}
    \end{block}

    \vspace{0.4em}
    \begin{alertblock}{Mise en cache}
      Rapport sauvegardé en JSON — pas de double appel LLM.
    \end{alertblock}
  \end{columns}

\end{frame}


% =============================================================================
\section{Conclusion}
% =============================================================================

\begin{frame}{Conclusion}

  \begin{columns}[T]
    \column{0.55\textwidth}
    \begin{block}{Ce que Persona accomplit}
      \begin{itemize}
        \item \textbf{15 questions} adaptatives, couverture OCEAN optimale
        \item Scoring Big Five normalisé
        \item Archétype par distance euclidienne
        \item Rapport narratif généré par \textbf{claude-opus-4-6}
        \item Transmission automatique par \textbf{e-mail}
      \end{itemize}
    \end{block}

    \column{0.42\textwidth}
    \begin{block}{Évolutions possibles}
      \begin{itemize}
        \item Modèle probabiliste (TRI)
        \item Déploiement cloud
        \item Historique multi-sessions
      \end{itemize}
    \end{block}
  \end{columns}

  \vspace{1.5em}
  \begin{center}
    {\Large\textbf{Merci pour votre attention}}\\[0.4em]
    \textcolor{lavender}{Questions ?}
  \end{center}

\end{frame}

\end{document}
